\subsection{Variables categóricas y abandono}

Se calculó el porcentaje de abandono respecto al total de clientes de cada categoría. Esta 
comparación evita confundir una categoría numerosa con una de mayor 
proporción de abandono. La Tabla~\ref{tab:categorias} reúne contrato, internet y método de pago.

\begin{table}[t]
    \caption{Clientes y abandono por categorías seleccionadas}
    \label{tab:categorias}\centering\footnotesize
    \begin{tabularx}{\columnwidth}{@{}Xrr@{}}\toprule
        Categoría & Clientes & Abandono (\%)\\\midrule
        \multicolumn{3}{@{}l}{\textit{Contrato}}\\
        Mensual & 3875 & 42.71\\
        Un año & 1473 & 11.27\\
        Dos años & 1695 & 2.83\\\midrule
        \multicolumn{3}{@{}l}{\textit{Internet}}\\
        DSL & 2421 & 18.96\\
        Fibra óptica & 3096 & 41.89\\
        Sin internet & 1526 & 7.40\\\midrule
        \multicolumn{3}{@{}l}{\textit{Método de pago}}\\
        Cheque electrónico & 2365 & 45.29\\
        Cheque por correo & 1612 & 19.11\\
        Transferencia automática & 1544 & 16.71\\
        Tarjeta de crédito automática & 1522 & 15.24\\\bottomrule
    \end{tabularx}
\end{table}

Los contratos mensuales presentan 42.71\% de abandono, frente a 11.27\% y 2.83\% en contratos de uno 
y dos años. La diferencia entre los extremos es de 39.88 puntos porcentuales. Esta diferencia muestra una 
asociación con el abandono, aunque los grupos también pueden diferir en antigüedad y servicios. Por eso, 
cambiar el contrato no garantiza que un cliente permanezca más tiempo.

La fibra óptica presenta 41.89\%, frente a 18.96\% en DSL. Como no se midió la velocidad, las 
fallas ni la satisfacción, estos resultados no permiten concluir que la fibra sea de peor calidad.
El cheque electrónico presenta 45.29\%, mientras que los otros métodos se sitúan entre 15.24\% y 
19.11\%. El método de pago podría estar relacionado con el perfil o contrato de los clientes, pero 
los datos no permiten confirmar la razón de esta diferencia.

Entre clientes con internet, la ausencia de seguridad en línea se asocia con 41.8\% de abandono, frente a 
14.6\% cuando se dispone de ella; para soporte técnico, los valores son 41.6\% y 15.2\%. La categoría 
\emph{No internet service} representa una situación distinta de no contratar el complemento y reúne a los 
mismos 1526 clientes en varias columnas. No debe contarse como evidencia independiente en cada comparación.

No todos los complementos presentan diferencias similares. Respaldo en línea muestra 39.9\% sin el 
complemento y 21.5\% con él; protección del dispositivo, 39.1\% y 22.5\%. Las diferencias son menores en 
streaming de televisión (33.5\% y 30.1\%) y películas (33.7\% y 29.9\%). Estos porcentajes comparan las 
categorías No y Yes y excluyen a quienes no tienen internet.

En género, los porcentajes son cercanos (26.9\% en Female y 26.2\% en Male). En \var{SeniorCitizen}, 
la categoría 1 presenta 41.7\%, frente a 23.6\% en la categoría 0. También hay diferencias por pareja 
(33.0\% sin pareja y 19.7\% con pareja), dependientes (31.3\% y 15.5\%) y facturación sin papel 
(33.6\% con ella y 16.3\% sin ella). Estas diferencias describen los grupos de la muestra. 
Todavía falta evaluar cuánto aportan estas variables al modelo y si sus errores varían entre grupos.
